\section{はじめに}\label{sec:introduction}

トランザクションデータにおけるアイテムセットマイニングは、小売分析、ウェブログ解析、医療データ分析など幅広い応用を持つ~\cite{agrawal1994fast,han2004mining}。
近年、単なる頻出パターンの発見にとどまらず、特定のアイテムセットが\emph{時間的に集中して出現する区間}（密集区間）を検出する手法が注目されている~\cite{suzuki2024dense}。
スライディングウィンドウを用いた密集区間検出は、サポート値の時間的変動を捉え、マーケティングキャンペーンの効果検証や異常検知に有用である。

しかし、トランザクションデータには個人の購買履歴や行動ログが含まれるため、マイニング結果から個人情報が推論されるリスクがある。
例えば、特定の薬品の組み合わせが密集する時間区間が公開されると、その期間に購入した個人の疾病状態が推測される可能性がある。
欧州一般データ保護規則（GDPR）や米国カリフォルニア州消費者プライバシー法（CCPA）の施行により、プライバシーを保証したデータマイニング手法の必要性は急速に高まっている~\cite{lee2011how}。

\emph{差分プライバシー}（Differential Privacy, DP）~\cite{dwork2006calibrating,dwork2014algorithmic}は、データベースへの1レコードの追加・削除が出力分布に与える影響を数学的に制限するフレームワークであり、プライバシー保証の黄金基準とされる。
頻出アイテムセットマイニングへの差分プライバシーの適用は既に研究されているが~\cite{zeng2012differentially,li2012privbasis,bhaskar2010discovering}、これらは全期間にわたるグローバルなサポート値を対象としている。
\emph{密集区間検出}では、各ウィンドウ位置におけるローカルなカウント値の系列に対して独立にプライバシー保護を行う必要があり、クエリ数のスケーリングや合成定理の適用において本質的に異なる課題が生じる。

密集区間マイニングに差分プライバシーを導入する際の主要な技術的課題は以下の通りである：
\begin{enumerate}
  \item \textbf{窓感度の解析}：スライディングウィンドウ内のカウントクエリに対する感度（1レコード変更の最大影響量）を正確に評価する必要がある。特に、1つのトランザクションが最大$W+1$個のウィンドウに影響するため、ベクトルクエリとしての感度と個別クエリとしての感度を区別する必要がある。
  \item \textbf{閾値判定の安定性}：ノイズ付加後のカウントに基づく密集/非密集の判定が、ノイズにより頻繁に反転する問題を制御する必要がある。
  \item \textbf{予算管理}：Apriori型の多レベル探索では多数のクエリが発行されるため、プライバシ予算の効率的な配分と合成が不可欠である。
\end{enumerate}

本研究の主要な貢献は以下の通りである：
\begin{itemize}
  \item ウィンドウカウントクエリの窓感度が$\Delta_W = 1$であることの形式的証明（定理\ref{thm:window-sensitivity}）。有界隣接関係に基づくプライバシー定義を明確化し、個別クエリ感度とベクトルクエリ感度の関係を論じる。
  \item ラプラスメカニズムおよびガウスメカニズムに基づくDP密集区間検出アルゴリズムDP-DIM（アルゴリズム\ref{alg:dp-dim}）と、そのプライバシー保証の正式な証明。
  \item 閾値安定性条件の導出と、安定マージンに基づく信頼度保証（定理\ref{thm:stability}）。
  \item Sparse Vector Technique を用いた予算効率的な変種（アルゴリズム\ref{alg:svt-dim}）。
  \item 多レベルApriori探索における階層的予算配分のプライバシー証明と、高度合成定理による予算削減の定量的分析。
  \item 合成データに対する包括的な実験評価（プライバシー・精度トレードオフ、メカニズム比較、パラメータ感度分析を含む）。
\end{itemize}

本論文の構成は以下の通りである。第\ref{sec:related}節で関連研究を概観し、第\ref{sec:preliminaries}節で予備知識を定義する。第\ref{sec:method}節で提案手法DP-DIMを詳述し、第\ref{sec:experiments}節で実験結果を示す。第\ref{sec:conclusion}節で結論と今後の展望を述べる。
